%%% ------------------------------------------------------------------
%%% An introduction to LaTeX for linguists.
%%%
%%% Adapted, with permission of the licence and with thanks, from
%%%   András Bárány, "Using LaTeX (in linguistics)", handout for the
%%%   2026 SGSAH summer school, Glasgow, 24 June 2026.
%%%   https://andras.barany.at/files/sgsah-latex-intro-handout.pdf
%%% This deck covers the first four pages of that handout (sections 1-3),
%%% with the exercises pointed at our own template instead of his.
%%% ------------------------------------------------------------------

\documentclass[a4paper,landscape,headrule,footrule]{foils}

%%
%%% macros for Theories of Grammar
%%%
\usepackage{polyglossia}
\setdefaultlanguage{english}
%\setmainfont{TeX Gyre Pagella}


\newcommand{\logo}{~}
\newcommand{\header}[3]{%
\title{\vspace*{-2ex} \large JMORF/MORS --- Morpho-Syntax
\\[2ex] \Large  \emp{#2} \\ \emp{#3}}
\author{\blu{Francis Bond}   \\ 
\normalsize  \textbf{Palacký University}\\
\normalsize  \url{https://fcbond.github.io/}\\
\normalsize  \texttt{bond@ieee.org}}
\MyLogo{JMORF (2026)}
\renewcommand{\logo}{#2}
\hypersetup{
   pdfinfo={
     Author={Francis Bond},
     Title={#1: #2},
     Subject={JMORF/MORS: Morpho-Syntax},
     Keywords={Syntax, Morphology, HPSG, Unification, Constructions},
     License={CC BY 4.0}
   }
 }

\date{#1
  \\ Location: KB-1.11}
}
\usepackage{hyperref}
\hypersetup{
     colorlinks,
     linkcolor={blue!50!black},
     citecolor={red!50!black},
     urlcolor={blue!80!black}
   }



\usepackage{xcolor}
\usepackage{graphicx}
\newcommand{\blu}[1]{\textcolor{blue}{#1}}
\newcommand{\grn}[1]{\textcolor{green}{#1}}
\newcommand{\hide}[1]{\textcolor{white}{#1}}
\newcommand{\emp}[1]{\textcolor{red}{#1}}
\newcommand{\txx}[1]{\textbf{\textcolor{blue}{#1}}}
\newcommand{\lex}[1]{\textbf{\mtcitestyle{#1}}}


% \renewcommand{\labelitemi}{\textcolor{violet}{‣}}
% \renewcommand{\labelitemii}{\textcolor{purple}{‣}}

\newcommand{\subhead}[1]{\noindent\textbf{#1}\\[5mm]}

\newcommand{\Bad}{\emp{\raisebox{0.15ex}{\ensuremath{\mathbf{\otimes}}}}}
\newcommand{\bad}[1]{*\eng{#1}}

\newcommand{\com}[1]{\hfill (\emp{#1})}%

\usepackage{relsize,xspace}
\newcommand{\into}{\ensuremath{\rightarrow}\xspace}
\newcommand{\tot}{\ensuremath{\leftrightarrow}\xspace}
\usepackage{url}
\newcommand{\lurl}[1]{\MyLogo{\url{#1}}}

\usepackage{langsci-gb4e}
\let\eachwordone=\slshape
%%% langsci-gb4e defines \thexnumi as \@roman\@xsi{xnumi}, which only works
%%% inside an exe environment; some decks use xlisti on its own.  Restore
%%% the plain gb4e definition so that keeps working.
\makeatletter
\def\thexnumi{\@xsi{xnumi}}
\makeatother
\newcommand{\lx}[1]{\textbf{\mtciteform{#1}}}
\newcommand{\ix}{\ex\slshape}



\newcommand{\ent}{\ensuremath{\Rightarrow}\xspace}
\newcommand{\ngv}{\ensuremath{\not\Rightarrow}\xspace}
%\usepackage{times}
%\usepackage{nttfoilhead}
\newcommand{\myslide}[1]{\foilhead[-25mm]{\raisebox{12mm}[0mm]{\emp{#1}}}\MyLogo{\logo}}
\newcommand{\myslider}[1]{\rotatefoilhead[-25mm]{\raisebox{12mm}[0mm]{\emp{#1}}}}
%\newcommand{\myslider}[1]{\rotatefoilhead{\raisebox{-8mm}{\emp{#1}}}}

\newcommand{\section}[1]{\myslide{}{\begin{center}\Huge \emp{#1}\end{center}}}



\usepackage[lyons,j,e,k]{mtg2e}
%\renewcommand{\mtcitestyle}[1]{\textcolor{-red!75!green!50}{\textsl{#1}}}
\renewcommand{\mtcitestyle}[1]{\textcolor{teal}{\textsl{#1}}}
\newcommand{\cs}{\mtciteform}
\newcommand{\iz}[1]{\texttt{\textup{#1}}}
\newcommand{\gm}{\textsc}
\usepackage[normalem]{ulem}
\newcommand{\ul}{\uline}
\newcommand{\ull}{\uuline}
\newcommand{\wl}{\uwave}
\newcommand{\vs}{\ensuremath{\Leftrightarrow}~}
%%%
%%% Bibliography
%%%
\usepackage{natbib}
%\usepackage{url}
\usepackage{bibentry}



\usepackage[utf8]{inputenc}
%%% trees
\usepackage{multicol}
%%% pst-node: one example in lec-12-control draws a connector between
%%% two points with \rnode/\ncbar.  This used to come in with rtrees.
\usepackage{pst-node}
\usepackage{forest}
\useforestlibrary{linguistics}
\forestapplylibrarydefaults{linguistics}

%%% \makebox's optional arguments are read as AVM delimiters by
%%% langsci-avm, so hide them behind a macro with no visible brackets.
\newcommand{\lbox}[1]{\makebox[0mm][l]{#1}}

%%% Tree styles.
%%%   words  -- trees whose leaves are actual words (italicised, aligned)
%%%   cat    -- trees whose leaves are categories or feature structures
\forestset{
  words/.style={for tree={s sep=7mm, l sep=5mm,
                          if n children=0{font=\itshape, tier=terminal}{}}},
  cat/.style={for tree={s sep=7mm, l sep=5mm}},
}
\newcommand{\wtree}[1]{\begin{forest} words #1 \end{forest}}
\newcommand{\ctree}[1]{\begin{forest} cat #1 \end{forest}}
%%% the rule being applied at this step of a derivation
\newcommand{\step}[1]{\begin{center}\emp{#1}\end{center}}

%%% AVMs (Language Science Press)
\usepackage{langsci-avm}

%%% Macros the decks use, carried over from headx.tex
\newcommand{\blank}{\rule{3em}{1pt}\xspace}
\newcommand{\ft}[1]{\textsc{#1}}
\newcommand{\idn}{\mtciteform}    %%% Indonesian
\newcommand{\prd}[1]{\textbf{#1}}
\newcommand{\typ}[1]{\textit{#1}}
%%% \val: provide before renewing, in case a package defines it.
\providecommand{\val}[1]{}
\renewcommand{\val}[1]{\textit{#1}}
%%% Manning's \avmbox, still used in tree labels and in maths where a
%%% bare langsci \tag would not work.  \avm{\tag{...}} gives the same
%%% boxed tag as the ones inside feature structures.
\newcommand{\avmbox}[1]{\mbox{\avm{\tag{#1}}}}

%%% Ring a value in a dotted oval.  These used pstricks \ovalnode,
%%% which came in with rtrees; tikz (via forest) does the same job.
\newcommand{\OV}[1]{\tikz[baseline]{\node[draw,dotted,red,ellipse,
          inner sep=1pt,outer sep=0pt]{#1};}}
\newcommand{\OVB}[1]{\tikz[baseline]{\node[draw,dotted,blue,ellipse,
          inner sep=1pt,outer sep=0pt]{#1};}}

%%% From CSLI book
\newcommand{\mc}{\multicolumn}
\newcommand{\HD}{\textbf{H}\xspace}
\newcommand{\el}{< >}       %%% empty list; langsci-avm needs the space

\newcommand{\A}{\noindent\textbf{A}: }
\newcommand{\Q}{\noindent\textbf{Q}: }
%\newcommand{\C}{\noindent\textbf{C}: }



%%% Hangul and IPA are not in the default font
\newfontfamily\korfont{Noto Sans CJK KR}
\newcommand{\ko}[1]{{\korfont #1}}
\newfontfamily\ipafont{Noto Sans}
\newcommand{\ipa}[1]{{\ipafont #1}}

%%% code listings
\usepackage{listings}
\lstset{
  language=[LaTeX]TeX,
  basicstyle=\ttfamily\footnotesize,
  keywordstyle=\color{blue!60!black}\bfseries,
  commentstyle=\color{teal}\itshape,
  stringstyle=\color{black},
  showstringspaces=false,
  columns=fullflexible,
  frame=single,
  framerule=0pt,
  backgroundcolor=\color{black!5},
  xleftmargin=1em,
  numbers=left,
  numberstyle=\tiny\color{black!50},
  numbersep=8pt,
}

%%% No font here has both Hangul and IPA, and listings will not let us
%%% switch font mid-listing, so the Korean source is set with alltt --
%%% verbatim-like, but commands still work, so \ko and \ipa can be used.
\usepackage{alltt}
\newfontfamily\monouni{DejaVu Sans Mono}
%%% An environment, not a command: alltt changes catcodes, so it cannot
%%% be passed as a macro argument -- the line breaks would be eaten.
%%% lrbox captures the content as it is typeset, which is fine.
\newsavebox{\srcb}
\newenvironment{srcbox}
  {\begin{lrbox}{\srcb}\begin{minipage}{0.93\linewidth}\monouni\footnotesize}
  {\end{minipage}\end{lrbox}%
   \noindent\colorbox{black!5}{\usebox{\srcb}}}
\newcommand{\bs}{\textbackslash}

\begin{document}
\header{LaTeX}{Using LaTeX in linguistics}{}
\maketitle

\myslide{Overview}
\begin{itemize}
\item What LaTeX is, and what it is not
\item Why you might want it
\item What it is good at, for us in particular
\item The basics: a document, commands, environments, formatting
\item Next week we will test it, so make an account on \href{https://overleaf.com}{Overleaf}
\end{itemize}

\bigskip
These slides follow the first part of a \href{https://andras.barany.at/files/sgsah-latex-intro-handout.pdf}{handout} by \emp{András Bárány}
--- see the acknowledgments at the end.

%%% ------------------------------------------------------------------
\section{What is LaTeX?}
%%% ------------------------------------------------------------------

\myslide{What is LaTeX?}

LaTeX is a \txx{typesetting system} for making documents.

\begin{itemize}
\item Scientific papers and presentations
\item Your thesis
\item Letters, reports, books, \ldots
\end{itemize}

\bigskip
For this course: your problem sets, and your assignment.

\myslide{What LaTeX is \emph{not}}

LaTeX is \emp{not} \txx{what you see is what you get}, unlike Word or
LibreOffice.

\bigskip
In a word processor, the thing you edit \emph{is} the thing you get.
In LaTeX the document is split up:

\begin{itemize}
\item You edit a plain text file ending in \iz{.tex}
\item That file is \txx{rendered} to give an output file
\item The output is normally a \iz{.pdf}
\item Other files can be pulled in: images, bibliographies, \ldots
\item You can output  \iz{.tex} from a program, and \txx{include} it in your document
  
\end{itemize}

\myslide{Why would you want that?}

\begin{itemize}
\item Separating the content from the output gives you \txx{flexibility}
\item The same content can be rendered differently
\item You can make slides \emph{and} a handout from one source
\item Editing a text file is \txx{more precise}: you say exactly what
  you mean, and it does not change under you
\item It works well with version control systems (like github)
\end{itemize}

\bigskip
It is also \txx{stable}: a file from twenty years ago still builds, and
builds the same.

%%% ------------------------------------------------------------------
\section{What it is good at}
%%% ------------------------------------------------------------------

\myslide{What LaTeX can do for a linguist}

\begin{itemize}
\item Bibliographies that look after themselves (\iz{biblatex}, \iz{natbib})
\item Long documents with many chapters
\item Contents, lists of figures, lists of abbreviations
\item Unicode: IPA, Hangul, Chinese, \ldots
\item Numbered examples, glosses, syntax trees, AVMs
\end{itemize}

\bigskip
The last one is why we use it in this course.

\myslide{Unicode: it just works}
\MyLogo{Although showing the code was complicated}

\begin{exe}
\ex Korean \\
  \ko{안녕하세요} \ipa{[an̠ɲʌŋɦasʰejo]} ``hello'' [kor]
\end{exe}

\begin{srcbox}\begin{alltt}
\bs{}begin\{exe\}
\bs{}ex Korean \bs\bs
  \ko{안녕하세요} \ipa{[an̠ɲʌŋɦasʰejo]} ``hello'' [kor]
\bs{}end\{exe\}
\end{alltt}\end{srcbox}

You type the characters and they come out --- as long as you compile
with \iz{xelatex} and pick a font that has them.

\myslide{Glossed examples}

\begin{exe}
\ex Hungarian [hun] \\
  \gll Ez  egy  komplex-ebb  példa. \\
       \gm{dem} \gm{indf.det} complex-\gm{cmpr} example \\
  \trans ``This is a more complex example.''
\end{exe}

\begin{lstlisting}
\begin{exe}
\ex Hungarian [hun] \\
  \gll Ez  egy  komplex-ebb  példa. \\
       \textsc{dem} \textsc{indf.det} complex-\textsc{cmpr} example \\
  \trans ``This is a more complex example.''
\end{exe}
\end{lstlisting}

The two lines line up by themselves, and the example is numbered for
you --- inserting one above does not mean renumbering by hand.

%%% ------------------------------------------------------------------
\section{The basics}
%%% ------------------------------------------------------------------

\myslide{A simple document}

A \iz{.tex} file needs at least two things:

\begin{itemize}
\item a \txx{preamble}, which sets up the form of the output
\item a \txx{document} environment, which holds the content
\end{itemize}

\begin{lstlisting}
\documentclass{article}

\begin{document}
Hello, world.
% This is a comment.
\end{document}
\end{lstlisting}

\myslide{Commands and environments}

LaTeX has two kinds of thing.

\begin{itemize}
\item A \txx{command} starts with a backslash, then a name, then zero or
  more arguments in braces
  \\ \iz{\textbackslash textbf\{bold\}} gives \textbf{bold}
\item An \txx{environment} is opened with
  \iz{\textbackslash begin\{\ldots\}} and closed with
  \iz{\textbackslash end\{\ldots\}}
  \\ this list is an \iz{itemize} environment
  \\ the content of a file sits in a \iz{document} environment
\end{itemize}

\bigskip
Anything after a \iz{\%} is a \txx{comment} and is ignored.

\myslide{Simple formatting}

\begin{itemize}
\item \iz{\textbackslash emph\{\ldots\}} for \emph{emphasis}
\item \iz{\textbackslash textbf\{\ldots\}} for \textbf{bold}
\item \iz{\textbackslash textsc\{\ldots\}} for \textsc{small capitals}
\item \iz{\textbackslash texttt\{\ldots\}} for \texttt{typewritten text}
\item \iz{\textbackslash footnote\{\ldots\}} for a footnote
\end{itemize}

\begin{lstlisting}
\begin{itemize}
  \item You \emph{emphasise} text with \emph{...}
  \item Use \textbf{...} for bold text
  \item Use \textsc{...} for small capitals
\end{itemize}
\end{lstlisting}

\myslide{Titles and sections}

\begin{itemize}
\item Sections and subsections are \txx{numbered for you}
\item So are examples, figures, tables and footnotes
\item A title is made from a title, an author and a date
\end{itemize}

\begin{lstlisting}
\title{An interesting phenomenon}
\author{Kim Smith \and Sandy Jones}
\date{31 January 2027}
\maketitle

\section{Introduction}
\end{lstlisting}

\myslide{Next week}

You will write your first document next week.

\begin{itemize}
\item Make an account at \url{https://www.overleaf.com/}
\item Start a new project and upload our template:
  \\ \iz{assignment.tex} and \iz{assignment.bib} from the
  \href{https://bond-lab.github.io/Theories-of-Grammar/assignment.html}{assignment
    page}
\item Press \emp{Recompile} and see what comes out
\item Then change the title to your own and recompile again
\end{itemize}

\bigskip
If something breaks, read the \emph{first} error message: the ones after
it are usually only consequences of the first.

\myslide{Acknowledgments and References}

\begin{itemize}
\item These slides follow the first four pages of\\
  \emp{András Bárány} (2026) \textit{Using LaTeX (in linguistics)},
  handout for the SGSAH summer school, Glasgow, 24 June 2026.
  \\ \url{https://andras.barany.at/files/sgsah-latex-intro-handout.pdf}
  \\ Used with thanks; the examples and the order of presentation are his, I have tweaked it a little and   the exercises have been pointed at our own template.
\item For the linguistics-specific packages see my
  \href{https://www.overleaf.com/read/hnjksbxgbhtq}{short guide to
    formatting HPSG with LaTeX}
\item \href{https://www.overleaf.com/learn}{Overleaf's own documentation}
  is good, and searchable
\end{itemize}

\end{document}

%%% Local Variables:
%%% coding: utf-8
%%% mode: latex
%%% TeX-PDF-mode: t
%%% TeX-engine: xetex
%%% End:
